\documentclass{article}
\usepackage[T1]{fontenc}
\usepackage{lmodern}
\usepackage{graphicx}
\usepackage{amsmath,amssymb}
\usepackage[a4paper,margin=2cm]{geometry}
\usepackage[absolute,overlay]{textpos}
\setlength{\TPHorizModule}{1mm}
\setlength{\TPVertModule}{1mm}
\usepackage[indonesian]{babel}
\usepackage{hyperref}
\setlength{\emergencystretch}{2em}
% unit: Halaman 4

\begin{document}

% === Halaman 4 (llm) ===

\section*{Teori Informasi (Information Theory)}

\begin{itemize}
    \item \textit{Information theory}: cabang ilmu yang mempelajari \textbf{kuantisasi, penyimpanan, transmisi, dan pengolahan informasi}
    \item Contoh kuantisasi:
    \begin{itemize}
        \item 1 bit untuk mengkodekan jenis kelamin (M/F)
        \item 3 bit untuk mengkodekan nama hari (ada 7 hari)
        \item 4 bit untuk mengkodekan 0 s/d 9
    \end{itemize}
    \item Contoh penyimpanan: pemampatan data untuk mengurangi ukuran ruang \textit{storage}
\end{itemize}

\end{document}
